\section{Introduction}
\label{sec:intro}

Safety-critical \CPS platforms increasingly pair a high-performance
application complex with a smaller supervisory core that can monitor the
main control stack and trigger a fail-safe response when needed. This
pattern appears in current automotive platforms and fits naturally with
functional-safety practice~\cite{ISO26262,ArmAutoSolutionsOverview,%
NVIDIADriveFSI}. In parallel, runtime monitoring techniques such as
\STL provide mature ways to express bounded-window requirements over
realized behavior~\cite{MalerNickovic2004,FainekosPappas2009,%
Deshmukh2017,Bartocci2018CPSMonitoring}. What remains less explicit is
the communication substrate between the controller and the monitor. That
substrate belongs squarely to computer architecture: it determines how
producer updates interact with private caches, ordering rules, and
interference from the rest of the chip.

That substrate matters because the monitor consumes a \emph{record}, not
a scalar. A practical control record contains an epoch, a timestamp,
multiple actuation fields, and metadata. Even on a cache-coherent
multicore, the monitor can still observe a mixed snapshot unless the
record is published with an explicit synchronization discipline. A policy
engine that reasons over such a torn record is evaluating a state the
system was never in.

Two design choices interact in that channel. The first is
\emph{transport}: the monitor can read the producer's record directly
through coherent shared memory, or it can read a mirror refreshed by an
explicit transfer path. The second is the \emph{publication primitive}
that binds field updates into one logical record. This paper separates
those choices explicitly. We compare five architectures that combine two
transports, direct coherent sharing and DMA transfer to a monitor-local
mirror, with three publication disciplines, none, a sequence lock, and
a generation-based double buffer. The evaluation uses a PWM-inspired
producer-monitor benchmark in gem5 ARM SE mode with Ruby
MESI\_Two\_Level.

We make four contributions.
\begin{itemize}
  \item \textbf{Coherence alone is not enough.} Both direct sharing and
    DMA mirroring produce torn snapshots when no publication primitive is
    used.
  \item \textbf{Primitive choice dominates transport choice.} Adding a
    DMA engine does not fix correctness by itself. In our matrix,
    \dmanaive tears more often than the unsynchronized direct-sharing
    baseline.
  \item \textbf{Generation-based double buffering is the strongest
    correct design in this study.} It achieves zero torn reads with the
    lowest CPU-side coherence traffic, the shortest whole-workload
    simulated time, and negligible reader retries.
  \item \textbf{Workload shape changes cost but not the ranking.} Under
    sustained bias, the direct-sharing protocols pay substantially more
    CPU-side traffic, but the ordering of the correct designs remains
    stable.
\end{itemize}

The paper is an architecture study, not a new detection method. The
monitor is included because it exercises the channel with a realistic
consumer that imposes a correctness contract and a timing-relevant
endpoint. The contribution is the characterization of the publication
path and the decomposition of cost into transport effects, primitive
effects, and read-side retry behavior.
